% Compatibilidad con instalaciones MiKTeX cuyo formato base es anterior a
% los paquetes Beamer/Hyperref instalados.
\providecommand{\IfFormatAtLeastTF}[3]{#3}
\providecommand{\IfFormatAtLeastT}[2]{}
\providecommand{\IfFormatAtLeastF}[2]{#2}
\providecommand{\IfDocumentMetadataTF}[2]{#2}
\providecommand{\IfDocumentMetadataT}[1]{}
\providecommand{\IfDocumentMetadataF}[1]{#1}
\makeatletter
\@namedef{ver@epstopdf-base.sty}{9999/12/31}
\makeatother

\documentclass[aspectratio=169,11pt]{beamer}

% =============================================================================
% MACROECONOMIA II | LICENCIATURA EN ECONOMIA | CIDE
% SESION 8: NUEVA ECONOMIA KEYNESIANA Y PRECIOS RIGIDOS
% =============================================================================

\usepackage[utf8]{inputenc}
\usepackage[english]{babel}
\uselanguage{English}
\languagepath{English}
\usepackage{graphicx}
\usepackage{amsmath,amssymb}
\IfFileExists{booktabs.sty}{
  \usepackage{booktabs}
}{
  \providecommand{\toprule}{\hline}
  \providecommand{\midrule}{\hline}
  \providecommand{\bottomrule}{\hline}
  \providecommand{\cmidrule}[1]{\hline}
}
\providecommand{\addlinespace}[1][0.5em]{}
\usepackage{xcolor}
\usepackage{ragged2e}
\usepackage{etoolbox}
\usepackage{array,tabularx}
\IfFileExists{microtype.sty}{\usepackage{microtype}}{}

\newtheorem{proposition}[theorem]{Proposición}
\newtheorem{remark}[theorem]{Observación}
\newtheorem{assumption}[theorem]{Supuesto}

\tolerance=1000
\emergencystretch=1em
\AtBeginDocument{\justifying}
\addtobeamertemplate{itemize body begin}{\justifying}{}
\addtobeamertemplate{enumerate body begin}{\justifying}{}
\addtobeamertemplate{description body begin}{\justifying}{}
\addtobeamertemplate{block begin}{}{\justifying}
\addtobeamertemplate{block alerted begin}{}{\justifying}
\addtobeamertemplate{block example begin}{}{\justifying}
\AtBeginEnvironment{column}{\justifying}

\definecolor{CIDENavy}{RGB}{3,25,54}
\definecolor{CIDEBlue}{RGB}{5,35,75}
\definecolor{CIDEInk}{RGB}{25,34,45}
\definecolor{CIDEGray}{RGB}{101,112,126}
\definecolor{CIDEGrayLight}{RGB}{243,244,246}

\setbeamercolor{normal text}{fg=CIDEInk,bg=white}
\setbeamercolor{structure}{fg=CIDEBlue}
\setbeamercolor{alerted text}{fg=CIDEBlue}
\setbeamercolor{title}{bg=CIDEBlue,fg=white}
\setbeamercolor{subtitle}{fg=CIDEInk}
\setbeamercolor{frametitle}{bg=CIDEBlue,fg=white}
\setbeamercolor{palette primary}{bg=CIDEBlue,fg=white}
\setbeamercolor{palette secondary}{bg=CIDEBlue,fg=white}
\setbeamercolor{palette tertiary}{bg=CIDENavy,fg=white}
\setbeamercolor{palette quaternary}{bg=CIDENavy,fg=white}
\setbeamercolor{section in head/foot}{bg=CIDENavy,fg=white}
\setbeamercolor{section in head/foot shaded}{bg=CIDENavy,fg=white!50}
\setbeamercolor{mini frame}{bg=white,fg=white}
\setbeamercolor{mini frame shaded}{bg=white!35,fg=white!35}
\setbeamercolor{upper separation line head}{bg=CIDENavy}
\setbeamercolor{lower separation line head}{bg=CIDEBlue}
\setbeamercolor{block title}{bg=CIDEBlue,fg=white}
\setbeamercolor{block body}{bg=CIDEGrayLight,fg=CIDEInk}
\setbeamercolor{block title alerted}{bg=CIDENavy,fg=white}
\setbeamercolor{block body alerted}{bg=CIDEGrayLight,fg=CIDEInk}
\setbeamercolor{caption name}{fg=CIDEBlue}

\setbeamerfont{author in head/foot}{size=\tiny}
\setbeamerfont{institute in head/foot}{size=\tiny}
\setbeamerfont{title in head/foot}{size=\tiny}
\setbeamersize{text margin left=0.70cm,text margin right=0.70cm}
\setlength{\leftmargini}{1.2em}
\setlength{\leftmarginii}{1.1em}
\setlength{\abovecaptionskip}{0.25em}
\setlength{\belowcaptionskip}{0pt}

\useoutertheme[subsection=false]{miniframes}
\setbeamertemplate{navigation symbols}{}
\setbeamertemplate{mini frame}[circle]
\setbeamertemplate{mini frame in current subsection}[circle]
\setbeamertemplate{section in head/foot}{\color{white}\insertsectionhead}
\setbeamertemplate{section in head/foot shaded}{\color{white!55}\insertsectionhead}
\setbeamertemplate{caption}[numbered]
\setbeamertemplate{section in toc}[circle]
\setbeamertemplate{subsection in toc}[circle]
\setbeamertemplate{itemize item}{\raisebox{0.15ex}{\tiny$\blacktriangleright$}}
\setbeamertemplate{itemize subitem}{\raisebox{0.15ex}{\tiny$\bullet$}}
\setbeamertemplate{enumerate item}[circle]

\setbeamertemplate{frametitle}{%
  \nointerlineskip
  \begin{beamercolorbox}[
    wd=\paperwidth,ht=3.05ex,dp=1.05ex,
    leftskip=0.42cm,rightskip=0.42cm
  ]{frametitle}
    \usebeamerfont{frametitle}\insertframetitle
  \end{beamercolorbox}%
}

\setbeamertemplate{footline}{%
  \leavevmode
  \vbox{%
    \offinterlineskip
    \hbox{%
      \begin{beamercolorbox}[wd=.72\paperwidth,ht=2.05ex,dp=.55ex,leftskip=.35cm]{palette secondary}
        \usebeamerfont{author in head/foot}\insertshortauthor
      \end{beamercolorbox}%
      \begin{beamercolorbox}[wd=.28\paperwidth,ht=2.05ex,dp=.55ex,rightskip=.35cm plus1fil]{palette secondary}
        \hfill\usebeamerfont{institute in head/foot}\insertshortinstitute
      \end{beamercolorbox}%
    }%
    \hbox{%
      \begin{beamercolorbox}[wd=.86\paperwidth,ht=1.85ex,dp=.50ex,leftskip=.35cm]{palette tertiary}
        \usebeamerfont{title in head/foot}\insertshorttitle
      \end{beamercolorbox}%
      \begin{beamercolorbox}[wd=.14\paperwidth,ht=1.85ex,dp=.50ex,center]{palette tertiary}
        \usebeamerfont{title in head/foot}\insertframenumber{} / \inserttotalframenumber
      \end{beamercolorbox}%
    }%
  }%
  \vskip0pt
}

\newcommand{\RutaLogoCIDE}{Figuras/General/logocide.png}
\setbeamertemplate{background}{%
  \vbox to \paperheight{%
    \vfill
    \hbox to \paperwidth{%
      \hfill
      \includegraphics[height=0.58cm,keepaspectratio]{\RutaLogoCIDE}%
      \hspace{0.18cm}%
    }%
    \vspace{0.68cm}%
  }%
}

\setbeamertemplate{title page}{%
  \vspace*{0.35cm}
  \centering
  \begin{beamercolorbox}[wd=\textwidth,sep=0.24cm,center,rounded=false]{title}
    \usebeamerfont{title}\inserttitle\par
  \end{beamercolorbox}
  \vspace{0.25cm}
  {\usebeamerfont{subtitle}\usebeamercolor[fg]{subtitle}\insertsubtitle\par}
  \vfill
  {\usebeamerfont{author}\insertauthor\par}
  \vspace{0.35cm}
  {\usebeamerfont{institute}\insertinstitute\par}
  \vfill
  {\usebeamerfont{date}\insertdate\par}
  \vspace{0.25cm}
}

\newcommand{\FigurasSesion}{Figuras/08_Nueva Economia Keynesiana}
\newcommand{\figura}[2][0.65\textheight]{%
  \includegraphics[width=\linewidth,height=#1,keepaspectratio]{\FigurasSesion/#2}%
}
\newcommand{\cidehighlight}[1]{\textcolor{CIDEBlue}{\bfseries #1}}
\newcommand{\fuentefigura}[1]{%
  \vspace{0.03cm}\par{\raggedright\tiny\textit{Fuente:} #1\par}%
}
\newcommand{\notafuentes}[1]{\note{[Sources] #1}}
\newcolumntype{Y}{>{\raggedright\arraybackslash}X}

\hypersetup{
  colorlinks=true,
  urlcolor=blue,
  linkcolor=.,
  citecolor=.,
  filecolor=.
}

\newif\ifsectiontoc
\sectiontocfalse
\AtBeginSection[]{%
  \ifsectiontoc
    \begin{frame}{Contenido}
      \tableofcontents[currentsection,hideothersubsections]
    \end{frame}
  \fi
}

\title[Macroeconomía II]{Sesión 8 \\ Nueva Economía Keynesiana}
\subtitle{Modelo Nuevo-Keynesiano de Rigidez Nominal}
\author[LECO]{Arturo López}
\institute[CIDE]{Centro de Investigación y Docencia Económicas (CIDE)}
\date[Otoño 2026]{Licenciatura en Economía, Otoño 2026}

\begin{document}

\begin{frame}
  \titlepage
\end{frame}

% =============================================================================
\section{Precios rígidos}
% =============================================================================

\begin{frame}{Objetivos}
  \begin{itemize}
    \setlength{\itemsep}{0.55em}
    \item Construir un modelo intertemporal con \cidehighlight{precios rígidos} en el corto plazo.
    \item Explicar por qué el dinero deja de ser neutral y describir el mecanismo de transmisión monetaria.
    \item Comparar la estabilización mediante política monetaria y política fiscal.
    \item Retomar la trampa de liquidez, la regla de Taylor y analizar las principales críticas al modelo.
  \end{itemize}
  \notafuentes{Williamson (2018), cap. 14, pp. 529--550.}
\end{frame}


\begin{frame}{Precios rígidos}
  \begin{itemize}
    \item Los Keynesianos sostienen que los salarios nominales y los precios son \cidehighlight{rígidos} (\textit{sticky}) en el \cidehighlight{corto plazo}.
    
    \item Por ello, la oferta y la demanda pueden no coincidir en todos los mercados en cada momento.
    
    \item Las fluctuaciones del ciclo pueden entonces reflejar un \cidehighlight{ajuste nominal incompleto}, generando brechas y asignaciones ineficientes.
    
    \item En este contexto, la política económica puede suavizar el ciclo y mejorar la asignación de recursos.
    
    \item Estos modelos son ampliamente utilizados en la investigación macroeconómica y en el diseño de política económica.
  \end{itemize}
\end{frame}



\begin{frame}{¿Por qué pueden ser rígidos los precios?}
  \small
  \begin{itemize}
    \setlength{\itemsep}{0.52em}
    \item \textbf{Costos de menú:} cambiar precios implica costos administrativos, de información o de relación con clientes (\href{https://www.jstor.org/stable/1885395}{Mankiw, 1985}).
    \item \textbf{Ajuste aleatorio:} sólo una fracción de las empresas puede revisar su precio en cada periodo (\cidehighlight{Calvo pricing}: \href{https://www.sciencedirect.com/science/article/abs/pii/0304393283900600?via\%3Dihub}{Calvo, 1983}).
  \end{itemize}
\end{frame}


\begin{frame}{Competencia imperfecta y rigidez de precios}
  \begin{itemize}
    \item Hasta ahora hemos supuesto que las empresas operan bajo competencia perfecta y toman los precios como dados.
    
    \item Ahora suponemos \cidehighlight{competencia monopolística}: cada empresa produce un bien diferenciado y fija su precio sujeta a la demanda de los hogares.
    
    \item La competencia monopolística genera poder de mercado. Sin embargo, la rigidez de precios impide que la empresa reajuste libremente su precio en cada periodo.
    
    \item Como el precio puede permanecer fijo durante varios periodos, la empresa lo determina de manera \textit{forward-looking}, considerando la evolución esperada de sus costos y de la demanda.
  \end{itemize}
\end{frame}

\begin{frame}{Supuestos del modelo}
  \small
  \begin{itemize}
    \setlength{\itemsep}{0.52em}
    \item El precio actual está predeterminado: \(P= P^*\).
    \item La inflación esperada es cero, por lo que la relación de Fisher implica \(R=r\).
    \item El banco central opera a través de fijar un objetivo operacional para la tasa de interés \(r^*\), pero controla directamente la oferta monetaria.
    \item La política monetaria determina directamente la demanda de producto:
      \[
        Y^d(r)=C^d(Y,r)+I^d(r)+G,
      \]
    \item Las empresas producen la cantidad demandada al precio fijo, aunque preferirían producir otra cantidad.
    \item Con \(P\) fijo, las empresas ajustan producción y empleo en lugar de reducir inmediatamente su precio.
  \end{itemize}
\end{frame}

\begin{frame}{Determinación del equilibrio de corto plazo}
  \small
  \begin{enumerate}
    \setlength{\itemsep}{0.58em}
    \item Antes del periodo presente, las empresas determinan el nivel general de precios \(P= P^*\).
    \item El banco central anuncia la tasa de política monetaria \(r^*\).
    \item La curva de demanda de producto determina
      \[
        Y^*=Y^d(r^*).
      \]
    \item Las empresas contratan el trabajo necesario para producir \(Y^*\):
      \[
        Y^*=zF(K,N^*).
      \]
    \item El salario real se ajusta para inducir a los hogares a ofrecer ese nivel de empleo.
    \item El banco central inyecta o retira dinero del mercado monetario para alcanzar \(r^*\):
      \[
        M^s= P^* L\!\left(Y^d(r^*),r^*\right).
      \]
  \end{enumerate}
\end{frame}

\begin{frame}{El modelo Nuevo Keynesiano}
  \centering
  \figura[6.0cm]{fig_14_1_modelo_nk.png}
  \fuentefigura{Williamson (2018), figura 14.1.}
\end{frame}


\begin{frame}{Tasa natural y brecha de producto}
  \small
  \begin{itemize}
    \setlength{\itemsep}{0.55em}
    \item La \cidehighlight{tasa natural de interés} \(r^m\) es la tasa que vacía el mercado de bienes y servicios:
      \[
        Y^d(r^m)=Y^s(r^m)=Y^m.
      \]
    \item \(Y^m\) es el nivel eficiente o natural de producción. La diferencia entre el nivel natural de producción y el nivel efectivo es la \cidehighlight{brecha de producto} o \cidehighlight{\textit{output gap}}: 
    \item       \[
        Y^m-Y
      \]
    \item Si el banco central fija \(r^*>r^m\), entonces \(Y^d(r^*)<Y^m\) y surge una brecha recesiva:
    \item Cambios en el \textit{output gap} reflejan entonces fluctuaciones de demanda entorno al nivel natural de producción.
  \end{itemize}
\end{frame}

% =============================================================================
\section{No-neutralidad del dinero}
% =============================================================================


\begin{frame}{Una expansión monetaria}
  \small
  Suponga que la economía parte de su equilibrio de largo plazo. El banco central reduce su meta de \(r_1\) a \(r_2\):
  \[
    r_2<r_1
    \quad\Longrightarrow\quad
    C^d\uparrow,\ I^d\uparrow
    \quad\Longrightarrow\quad
    Y^d\uparrow.
  \]

  Como \(P=P_1\) está fijo, las empresas responden aumentando producción y empleo:
  \[
    Y_2>Y_1,
    \qquad N_2>N_1.
  \]

  Para implementar \(r_2\), el banco central satisface la mayor demanda monetaria:
  \[
    M_2=P_1L(Y_2,r_2)>P_1L(Y_1,r_1)=M_1.
  \]
\end{frame}

\begin{frame}{No neutralidad del dinero}
  \centering
  \figura[6.0cm]{fig_14_2_no_neutralidad.png}
  \fuentefigura{Williamson (2018), figura 14.2.}
\end{frame}

\begin{frame}{Mecanismo Keynesiano de transmisión monetaria}
  \small
  \[
    r\downarrow
    \ \Longrightarrow\ 
    M^s\uparrow
    \ \Longrightarrow\ 
    C^d\uparrow,\ I^d\uparrow
    \ \Longrightarrow\ 
    Y^d\uparrow
    \ \Longrightarrow\ 
    Y\uparrow,\ N\uparrow.
  \]

  \begin{itemize}
    \setlength{\itemsep}{0.55em}
    \item La expansión tiene primero efectos en el mercado financiero y después en el gasto agregado.
    \item La menor tasa reduce la oferta laboral deseada por sustitución intertemporal; el salario real debe aumentar para inducir más empleo actual.
    \item En el modelo: \(Y\), \(N\), \(w\), \(C\), \(I\) y \(M\) aumentan.
  \end{itemize}
\end{frame}

\begin{frame}{Corto plazo y largo plazo}
  \small
  \renewcommand{\arraystretch}{1.35}
  \begin{tabularx}{\textwidth}{lYY}
    \toprule
      & \textbf{Corto plazo} & \textbf{Largo plazo} \\
    \midrule
    Nivel de precios
      & Predeterminado: \(P=P^*\).
      & Ajusta hasta vaciar los mercados. \\
    Producción y empleo
      & Responden a la demanda agregada.
      & Regresan a sus niveles naturales. \\
    Dinero
      & Tiene efectos reales.
      & Es neutral en el modelo básico. \\
    Política monetaria
      & Puede cerrar la brecha de producto.
      & Determina principalmente variables nominales. \\
    \bottomrule
  \end{tabularx}
\end{frame}

% =============================================================================
\section{Política de estabilización}
% =============================================================================

\begin{frame}{¿Por qué estabilizar la economía?}
  \small
  \begin{itemize}
    \setlength{\itemsep}{0.6em}
    \item Dada la falla en los mercados privados para vaciarse en el corto plazo, un choque puede dejar a la economía fuera de su equilibrio de largo plazo.
    \item En particular, un choque puede dejar al precio vigente por encima del nivel que vaciaría el mercado de bienes.
    \item La producción demandada queda por debajo de la producción eficiente: \(Y<Y^m\).
    \item La asignación no es Pareto eficiente porque bienes y trabajo no se intercambian aun cuando sería mutuamente beneficioso hacerlo.
    \item El ajuste espontáneo ocurre cuando los precios disminuyen, pero puede ser lento.
    \item La política de estabilización intenta cerrar la brecha antes de que termine el ajuste nominal.
  \end{itemize}
\end{frame}

\begin{frame}{Estabilización mediante política monetaria}
  \centering
  \figura[6.35cm]{fig_14_3_estabilizacion_monetaria.png}
  \fuentefigura{Williamson (2018), figura 14.3.}
\end{frame}



\begin{frame}{Estabilización mediante política fiscal}
  \centering
  \figura[6.25cm]{fig_14_4_estabilizacion_fiscal.png}
  \fuentefigura{Williamson (2018), figura 14.4.}
\end{frame}






\begin{frame}{Choque de productividad}
  \centering
  \figura[6.35cm]{fig_14_5_ptf_politica_optima.png}
  \fuentefigura{Williamson (2018), figura 14.5.}
\end{frame}

\begin{frame}{Política endógena e identificación}
  \small
  \begin{itemize}
    \setlength{\itemsep}{0.55em}
    \item Un aumento persistente de la PTF desplaza \(Y^d\) y \(Y^s\) hacia la derecha.
    \item Sin respuesta monetaria, la producción aumenta sólo a \(Y_3\) y queda una brecha \(Y_2-Y_3\).
    \item Si el banco central reduce \(r_1\) a \(r_2\), la producción alcanza el nivel eficiente \(Y_2\).
    \item Con esa respuesta óptima, \(Y\), \(C\), \(I\) y \(N\) se comportan como en un modelo RBC.
  \end{itemize}

  \begin{alertblock}{Equivalencia observacional}
    Los mismos datos reales pueden provenir de precios flexibles o de precios rígidos acompañados por una política monetaria óptima. La política endógena dificulta identificar el mecanismo estructural.
  \end{alertblock}
\end{frame}


% =============================================================================
\section{Trampa de liquidez}
% =============================================================================

\begin{frame}{Política monetaria en el ZLB}
  \small
  \begin{enumerate}
    \setlength{\itemsep}{0.58em}
    \item Supongamos que la economía se encuentra en el ZLB: $R=0$.
    \item La economía parte de \((Y_1,r_1)\) con brecha nula.
    \item Ocurre un choque negativo de demanda agregada.
    \item La tasa que cerraría la brecha es \(r_2<r_1\).
    \item La tasa permanece en \(r_1\) y la producción cae a \(Y_2\).
    \item Si una tasa nominal negativa fuera factible, la economía podría alcanzar su nivel de producto potencial.
  \end{enumerate}
\end{frame}


\begin{frame}{Trampa de liquidez}
  \centering
  \figura[6.0cm]{fig_14_6_trampa_liquidez.png}
  \fuentefigura{Williamson (2018), figura 14.6.}
\end{frame}



\begin{frame}{Regla de Taylor y el ZLB}
  \centering
  \figura[6.35cm]{fig_14_7_regla_taylor.png}
  \fuentefigura{Williamson (2018), figura 14.7.}
\end{frame}





\end{document}
